% nucapp_pet: positron emission tomography (Vol III Ch 5, fig:nucapp-pet)
% A fluorine-18 nucleus in a tumour emits a positron, which travels about 1 mm, meets an
% electron and annihilates: e+ + e- -> 2 gamma.  The pair is nearly at rest, so total momentum
% is about zero and the two photons (each m_e c^2 = 511 keV) leave back to back.  Both strike the
% detector ring at the same instant, and the decay lies on the line joining the two detectors.
% Geometry: ring of 24 detectors, inner face at radius 2.55; annihilation point (0.42,-0.2); the line of
% response at 35 degrees; each photon is drawn to radius 2.62 (inside a detector).
\begin{tikzpicture}
  % detector ring
  \foreach \a in {0,15,...,345}
    \draw[fill=cTeal!35,draw=cTeal,thin line,rotate=\a] (2.55,-0.18) rectangle (2.95,0.18);
  % patient and tumour
  \draw[fill=cMoon,draw=cSteel,thin line] (0,0) ellipse[x radius=1.25,y radius=1.7];
  \node[cSteel,font=\footnotesize] at (0,1.3) {patient};
  \fill[cAmber!70] (0.35,-0.25) circle[radius=0.32];
  \node[font=\footnotesize,anchor=west] at (0.62,-0.62) {tumour};
  \coordinate (P) at (0.42,-0.2);
  % line of response
  \draw[guide] ($(P)+(35:-2.75)$) -- ($(P)+(35:2.45)$);
  % the two photons
  \draw[-bookhead,thin line,photon] (P) -- ($(P)+(35:2.36)$);
  \draw[-bookhead,thin line,photon] (P) -- ($(P)+(215:2.82)$);
  \node[dot] at (P) {};
  \node[cAcc,anchor=west] at (3.15,1.75) {$\gamma$, \SI{511}{keV}};
  \node[cAcc,anchor=east] at (-2.75,-2.05) {$\gamma$, \SI{511}{keV}};
  % explanation
  \node[anchor=west,align=left] at (3.35,-0.35) {$e^++e^-\to2\gamma$\\back to back:\\momentum conserved};
  \node[cTeal,anchor=west] at (3.35,-1.65) {detector ring};
\end{tikzpicture}
